% Source: https://artofproblemsolving.com/community/c3h1486770_determining_number_from_remainders
% Stable identity: aops-topic-1486770-q01
Let the number be $N$. The conditions say
\[
N\equiv19\equiv2\pmod{17},\qquad N\equiv-21\equiv17\pmod{19}.
\]
From the second, $N=17+19t$. Then $17+19t\equiv2\pmod{17}$ becomes $2t\equiv2\pmod{17}$ (since $19\equiv2$ and $17\equiv0$), so $t\equiv1\pmod{17}$. With $t=1+17m$,
\[
N=17+19(1+17m)=36+323m .
\]
The three-digit values are $36+323=\boxed{359}$ and $36+646=\boxed{682}$ (the next one, $1005$, has four digits). Check: $359-19=340=17\cdot20$, $359+21=380=19\cdot20$; $682-19=663=17\cdot39$, $682+21=703=19\cdot37$.
